\documentclass{article}
\usepackage{fontspec} 
\usepackage{unicode-math}  % 先加载 unicode-math
\setmathfont{XITS Math}    % 再设置数学字体
%\usepackage[UTF8]{ctex} % 如果需要支持中文

\usepackage{anyfontsize}
\usepackage{xeCJK} % XeLaTeX 推荐 xeCJK，LuaLaTeX 也可用   
\usepackage{setspace}  % 添加这个包
\setstretch{1.2}      % 设置行距为1.5倍

\setCJKmainfont[
    Path = C:/USERS/11252/APPDATA/LOCAL/MICROSOFT/WINDOWS/FONTS/,  % 改用 Windows 字体目录
    UprightFont = {SourceHanSansCN-Light1.otf},  % 使用可变字体
    BoldFont = {SourceHanSansCN-Medium1.otf},
    LetterSpace = 2.0,  % 增加字间距
    WordSpace = 1.2,    % 增加词间距
    %BoldFeatures = {RawFeature={+wght=900}},
    %Scale=1
]{SourceHanSansCN}


\setmainfont[
    Path = C:/USERS/11252/APPDATA/LOCAL/MICROSOFT/WINDOWS/FONTS/,  % 改用 Windows 字体目录
    UprightFont = {SourceHanSansCN-Light1.otf},  % 使用可变字体
    BoldFont = {SourceHanSansCN-Medium1.otf},
    LetterSpace = 2.0,  % 增加字间距
    WordSpace = 1.2,    % 增加词间距
    %BoldFeatures = {RawFeature={+wght=900}},
    %Scale=1
]{SourceHanSansCN}

%\setCJKmainfont[
 %   Path = C:/USERS/11252/APPDATA/LOCAL/MICROSOFT/WINDOWS/FONTS/,
 %   UprightFont = {SourceHanSerifCN-Light-5.otf},
 %   BoldFont = {SourceHanSerifCN-Medium-6.otf},
%    Scale=1
%]{SourceHanSerifCN}

%\setmainfont{SourceHanSerif}  % 设置英文字体

% 处理冲突的命令
\let\oldbackepsilon\backepsilon
\let\backepsilon\relax
\let\oldeth\eth
\let\eth\relax
\let\olddigamma\digamma
\let\digamma\relax
\usepackage{float}

\usepackage[a4paper, left=0.1in, right=0.1in, top=0.05in, bottom=0.05in]{geometry} % 设置四边页边距为0.1英寸{geometry} % 设置页边距为0.5英寸
\usepackage{enumitem} % 用于自定义列表
\usepackage{amsmath}  % 数学公式支持
\usepackage{tikz}
\usetikzlibrary{arrows.meta, positioning}
\setlength{\parindent}{0pt} % 不缩进
\usepackage{etoolbox} % 用于列表操作
%调整类代码
\usepackage{comment}
\usepackage{tocloft} % 用于定制目录 样式
\usepackage{hyperref} %不需要目录盒子注释这
\usepackage{ifthen}
\usepackage{tikz}
\usetikzlibrary{arrows.meta}
\usepackage{titletoc}
\usepackage{graphicx}
\usepackage{float}

\usepackage{amssymb}  % 提供数学符号，包括\therefore
%目录设定
%快捷键 CMD + / 注释
%  \renewcommand{\cftdot}{} % 隐藏目录中的页码
%  \cftpagenumbersoff{section}
%  \cftpagenumbersoff{subsection}
%  \makeatletter
%  \renewcommand{\@pnumwidth}{0em}  % 设置页码宽度为0
%  \renewcommand{\@tocrmarg}{0em}   % 设置右边距为0
%  \makeatother
 

\usepackage{environ}

\newif\ifshowcontent  % 定义一个条件判断变量
\showcontenttrue    % 设置为 false，隐藏所有 question 环境的内容

\newcounter{question} % 题目计数器
\setcounter{question}{0}
\NewEnviron{question}{%
    \ifshowcontent
        \par\stepcounter{question}\textbf{\thequestion.} \BODY\par % 如果显示内容，则输出
    \fi
}

\NewEnviron{choices}{%
    \ifshowcontent
        \begin{enumerate}[label=\Alph*., leftmargin=*]
            \BODY
        \end{enumerate}\par
    \fi
}

\NewEnviron{correctanswer}{%
    \ifshowcontent
        \textbf{答案:}\BODY\par
    \fi
}



\newenvironment{explanation}
{\textbf{解析:}\par}
{\par}



% 新增考察方面命令
\newcommand{\checklist}[1]{
    \textbf{考察:} #1 \par % 在题目下方显示考察方面
    \addcontentsline{toc}{subsection}{题号\thequestion 考察: #1} % 将考察内容添加到目录
    \vspace{2em} % 添加1em的垂直空间
}

\graphicspath{{images/}} %统一


\begin{document}
 %\fontsize{22}{31}\selectfont
 \tableofcontents % 添加目录
\newpage
%\pagestyle{empty}




\section{考点1:存款服务和非存款负债的管理}
\setcounter{question}{0}


\begin{question}
    Which of the following describes a certificate of deposit (CD) that carries an interest rate that periodically adjusts upward?
\end{question}

\begin{choices}
    \item Liquidity CD
    \item Bump-up CD
    \item Step-up CD
    \item Index CD
\end{choices}

\begin{correctanswer}
    C
\end{correctanswer}

\begin{explanation}
    \textbf{A选项错误。}
    这个说法不准确：
    \begin{itemize}
        \item 流动性CD指可提前支取的CD
        \item 与利率调整无关
    \end{itemize}

    \textbf{B选项错误。}
    这个说法不准确：
    \begin{itemize}
        \item 加息CD利率可能随市场利率上升
        \item 但不是定期自动调整
    \end{itemize}

    \textbf{C选项正确。}
    这个说法正确：
    \begin{itemize}
        \item 递增CD利率会定期自动上调
        \item 是预先设定的调整机制
    \end{itemize}

    \textbf{D选项错误。}
    这个说法不准确：
    \begin{itemize}
        \item 指数CD利率与特定指数挂钩
        \item 不是定期自动调整
    \end{itemize}
\end{explanation}

\checklist{定期存款类型识别}

\begin{question}
    If the bank raises its offer rate on new deposits from 7 percent to 7.5 percent, Table above shows the marginal cost of this change: Change in total cost = \$50 million × 7.5 percent - \$25 million × 7 percent = \$3.75 million - \$1.75 million = \$2.00 million. The marginal cost rate:
\end{question}

\begin{choices}
    \item 9\%
    \item 8\%
    \item 7\%
    \item 6\%
\end{choices}

\begin{correctanswer}
    B
\end{correctanswer}

\begin{explanation}
    \textbf{A选项错误。}
    这个说法不准确：
    \begin{itemize}
        \item 9\%是75百万存款时的边际成本率
        \item 不是当前情况
    \end{itemize}

    \textbf{B选项正确。}
    这个说法正确：
    \begin{itemize}
        \item 边际成本 = 总成本变化
        \item = 新利率×新资金 - 旧利率×旧资金
        \item = \$50百万×7.5\% - \$25百万×7\%
        \item = \$3.75百万 - \$1.75百万 = \$2百万
        \item 边际成本率 = 总成本变化/额外资金
        \item = \$2百万/\$25百万 = 8\%
    \end{itemize}

    \textbf{C选项错误。}
    这个说法不准确：
    \begin{itemize}
        \item 7\%是原始利率
        \item 不是边际成本率
    \end{itemize}

    \textbf{D选项错误。}
    这个说法不准确：
    \begin{itemize}
        \item 计算结果错误
        \item 低估了边际成本率
    \end{itemize}
\end{explanation}

\checklist{边际成本率计算}


\begin{question}
    An analyst in a commercial bank is evaluating the available funds gap for the next month. After analysis, he found that 1) the qualified loan requests accepted is \$150 Million, 2) a major customer of the bank is going to withdraw \$135 Million, 3) the bank has plan to invest in bond markets for \$50 Million, 4) the current deposit is \$160 Million and \$100 Million is expected to come next month, 5) \$10 Million will be added to available funds gap for a cushion. Which of the following is the amount of funds that the bank needs to finance from non-deposit sources?
\end{question}

\begin{choices}
    \item 85 Million
    \item 75 Million
    \item 65 Million
    \item 335 Million
\end{choices}

\begin{correctanswer}
    A
\end{correctanswer}

\begin{explanation}
    \textbf{A选项正确。}
    这个说法正确：
    \begin{itemize}
        \item 可用资金缺口(AFG)计算：
        \item = 贷款需求 + 提款 + 投资计划 - 当前存款 - 预期存款 + 缓冲
        \item = \$150M + \$135M + \$50M - \$160M - \$100M + \$10M
        \item = \$85M
    \end{itemize}

    \textbf{B选项错误。}
    这个说法不准确：
    \begin{itemize}
        \item 计算结果错误
        \item 低估了资金缺口
    \end{itemize}

    \textbf{C选项错误。}
    这个说法不准确：
    \begin{itemize}
        \item 计算结果错误
        \item 低估了资金缺口
    \end{itemize}

    \textbf{D选项错误。}
    这个说法不准确：
    \begin{itemize}
        \item 计算结果错误
        \item 高估了资金缺口
    \end{itemize}
\end{explanation}

\checklist{可用资金缺口计算}


\begin{question}
    A deposit institution has accepted \$30 million qualified loan request. The management is considering issuing \$30 million negotiable CDs with current interest rate of 6\% to finance funds, the noninterest cost in the form of flotation cost in 0.2\%, and legal reserve requirements is 5\% of the total loan amount. What is the effective financing cost rate by issuing negotiable CDs?
\end{question}

\begin{choices}
    \item 6.2\%
    \item 6.53\%
    \item 6.6\%
    \item 5.8\%
\end{choices}

\begin{correctanswer}
    B
\end{correctanswer}

\begin{explanation}
    \textbf{A选项错误。}
    这个说法不准确：
    \begin{itemize}
        \item 仅考虑了利率和发行成本
        \item 未考虑法定准备金要求
    \end{itemize}

    \textbf{B选项正确。}
    这个说法正确：
    \begin{itemize}
        \item 有效融资成本率计算：
        \item = (利息成本 + 非利息成本)/(1 - 准备金率)
        \item = \$30M × (6\% + 0.2\%)/[\$30M × (1 - 5\%)]
        \item ≈ 6.53\%
    \end{itemize}

    \textbf{C选项错误。}
    这个说法不准确：
    \begin{itemize}
        \item 计算结果错误
        \item 高估了融资成本
    \end{itemize}

    \textbf{D选项错误。}
    这个说法不准确：
    \begin{itemize}
        \item 计算结果错误
        \item 低估了融资成本
    \end{itemize}
\end{explanation}

\checklist{有效融资成本率计算}


\begin{question}
    James, an analyst in liability management department, has to propose an appropriate non-deposit financing sources for the bank's coming \$50 million investment requirement next week. After investigate, he gets the following information and both ways require deposit insurance fee. Which of the following statements are most likely accurate?
\end{question}

\begin{choices}
    \item The effective cost rate of the Fed funds is 6.3\%
    \item The effective cost rate of the commercial paper is 7.2\%
    \item James should suggest the Fed Funds for financing due to the lower effective cost rate
    \item James should suggest the commercial paper for financing due to the lower effective cost rate
\end{choices}

\begin{correctanswer}
    C
\end{correctanswer}

\begin{explanation}
    \textbf{A选项错误。}
    这个说法不准确：
    \begin{itemize}
        \item 联邦基金有效成本率计算错误
        \item 实际应为6.5\%
    \end{itemize}

    \textbf{B选项错误。}
    这个说法不准确：
    \begin{itemize}
        \item 商业票据有效成本率计算错误
        \item 实际应为7.5\%
    \end{itemize}

    \textbf{C选项正确。}
    这个说法正确：
    \begin{itemize}
        \item 联邦基金成本率 = \$50M × (6\% + 0.25\%)/(\$50M - \$2M - \$50M × 0.003) ≈ 6.5\%
        \item 商业票据成本率 = \$50M × (7\% + 0.2\%)/(\$50M - \$2M - \$50M × 0.003) ≈ 7.5\%
        \item 联邦基金成本率更低
    \end{itemize}

    \textbf{D选项错误。}
    这个说法不准确：
    \begin{itemize}
        \item 商业票据成本率更高
        \item 不应选择商业票据融资
    \end{itemize}
\end{explanation}

\checklist{融资方式选择}


\begin{question}
    While smaller banks and thrift institutions usually rely most heavily on deposits for their funding needs, leading depository institutions around the globe have come to regard the nondeposit funds market as a key source of short-term money to meet both loan demand and unexpected cash emergencies. The most popular domestic source of borrowed reserves among depository institutions is the Federal funds market. Which of the following statement about the Federal funds market is most likely accurate?
\end{question}

\begin{choices}
    \item Specific collateral is needed for all Federal funds
    \item The Fed has set three levels of credit, primary, secondary and seasonal, to meet needs of different institutions
    \item The Federal funds market allows depository institutions short of reserves to meet their legal reserve requirements or to satisfy loan demand by tapping immediately usable funds from other institutions possessing temporarily idle funds
    \item The Federal funds market improves the liquidity of home mortgages by allowing troubled institutions to use the home mortgages they held in their portfolios as collateral for emergency loans
\end{choices}

\begin{correctanswer}
    C
\end{correctanswer}

\begin{explanation}
    \textbf{A选项错误。}
    这个说法不准确：
    \begin{itemize}
        \item 隔夜资金通常不需要抵押
        \item 联邦基金市场主要是无抵押交易
    \end{itemize}

    \textbf{B选项错误。}
    这个说法不准确：
    \begin{itemize}
        \item 这是美联储银行的特征
        \item 不是联邦基金市场的特征
    \end{itemize}

    \textbf{C选项正确。}
    这个说法正确：
    \begin{itemize}
        \item 联邦基金市场允许储备不足的机构
        \item 从拥有闲置资金的机构获取即时可用资金
        \item 用于满足法定准备金要求或贷款需求
    \end{itemize}

    \textbf{D选项错误。}
    这个说法不准确：
    \begin{itemize}
        \item 这是联邦住房贷款银行(FHLB)的特征
        \item 不是联邦基金市场的特征
    \end{itemize}
\end{explanation}

\checklist{联邦基金市场特征}


\begin{question}
    A small commercial bank wants to know the minimum rate of return it has to earn for the next year. Jack a treasurer from the investing department, starts from the total cost of capital, he found that next year, all coming funds is \$400 million, however, 10\% of them cannot be invested, and the current average interest rate is 5\%, and operating expense, is about \$10 million. According to the information above, please help Jack figure out the hurdle rate for investing.
\end{question}

\begin{choices}
    \item 7.5\%
    \item 8.3\%
    \item 2.5\%
    \item 5.5\%
\end{choices}

\begin{correctanswer}
    B
\end{correctanswer}

\begin{explanation}
    \textbf{A选项错误。}
    这个说法不准确：
    \begin{itemize}
        \item 计算结果错误
        \item 低估了最低回报率
    \end{itemize}

    \textbf{B选项正确。}
    这个说法正确：
    \begin{itemize}
        \item 最低回报率计算：
        \item = (利息成本 + 运营成本)/可投资资金
        \item = (\$400M × 5\% + \$10M)/[\$400M × (1 - 10\%)]
        \item = 8.3\%
    \end{itemize}

    \textbf{C选项错误。}
    这个说法不准确：
    \begin{itemize}
        \item 计算结果错误
        \item 严重低估了最低回报率
    \end{itemize}

    \textbf{D选项错误。}
    这个说法不准确：
    \begin{itemize}
        \item 计算结果错误
        \item 低估了最低回报率
    \end{itemize}
\end{explanation}

\checklist{最低回报率计算}

 
\begin{question}
    In a deposit institution, savings deposits, checking accounts, NOW accounts, Christmas Club accounts, time deposits, cashiers' checks, money orders, officers'checks, and any outstanding drafts normally are protected by federal insurance. About the deposit insurance, which of the following is the least accurate?
\end{question}

\begin{choices}
    \item Deposit insurance allows institutions to sell deposits at relatively low rates of interest
    \item Only members of Federal Deposit Insurance Corporation (FDIC) have deposit insurance
    \item Liquidity risk for institutions is eliminated by deposit insurance because customers are less panic about the loss of their funds in deposit institutions
    \item There is a limit on the deposit insurance, customers only get part of their deposit back if any crisis happened
\end{choices}

\begin{correctanswer}
    C
\end{correctanswer}

\begin{explanation}
    \textbf{A选项错误。}
    这个说法不准确：
    \begin{itemize}
        \item 该说法正确
        \item 存款保险确实允许机构以较低利率吸收存款
    \end{itemize}

    \textbf{B选项错误。}
    这个说法不准确：
    \begin{itemize}
        \item 该说法正确
        \item 只有FDIC成员才能获得存款保险
    \end{itemize}

    \textbf{C选项正确。}
    这个说法正确：
    \begin{itemize}
        \item 存款保险不能消除流动性风险
        \item 只能缓释流动性风险
        \item 即使有保险，客户仍可能因恐慌而提款
    \end{itemize}

    \textbf{D选项错误。}
    这个说法不准确：
    \begin{itemize}
        \item 该说法正确
        \item 存款保险确实有上限
    \end{itemize}
\end{explanation}

\checklist{存款保险特征}


\begin{question}
    Important indicators of management's effectiveness are whether or not funds deposited by the public have been raised at the lowest possible cost and whether sufficient deposits are available to fund all those loans and projects management wishes to pursue. Which of the following statements about deposit pricing methods are most appropriate?
    \begin{enumerate}
        \item Conditional pricing is sensitive to the types of customers each depository institution plans to serve and the cost that serving different types of depositors will present to the offering institution.
        \item Marginal revenue should be determined at the beginning when using marginal cost pricing method.
        \item Free-pricing may attract too much volatile funds when market rate is high, which actually adverse to deposit institutions.
        \item Relationship pricing promotes greater customer loyalty and makes the customer more sensitive to the prices posted on services offered by competing financial firms.
    \end{enumerate}
\end{question}

\begin{choices}
    \item I \& II
    \item II \& III
    \item I, II \& III
    \item I, II, III \& IV
\end{choices}

\begin{correctanswer}
    C
\end{correctanswer}

\begin{explanation}
    \textbf{A选项错误。}
    这个说法不准确：
    \begin{itemize}
        \item 遗漏了III的正确说法
        \item 自由定价确实可能吸引过多不稳定资金
    \end{itemize}

    \textbf{B选项错误。}
    这个说法不准确：
    \begin{itemize}
        \item 遗漏了I的正确说法
        \item 条件定价确实与客户类型相关
    \end{itemize}

    \textbf{C选项正确。}
    这个说法正确：
    \begin{itemize}
        \item I正确：条件定价与客户类型相关
        \item II正确：边际成本定价需要先确定边际收入
        \item III正确：自由定价可能吸引不稳定资金
        \item IV错误：关系定价使客户对竞争对手价格更不敏感
    \end{itemize}

    \textbf{D选项错误。}
    这个说法不准确：
    \begin{itemize}
        \item 包含了IV的错误说法
        \item 关系定价使客户对竞争对手价格更不敏感
    \end{itemize}
\end{explanation}

\checklist{存款定价方法分析}


\begin{question}
    Deposits provide much of the raw material for making loans and, thus, may represent the ultimate source of profits and growth for a depository institution. A deposit manager is trying to make the appropriate interest rate for deposits in order to balance the cost and profit. Among the following statements, which is correct?
\end{question}

\begin{choices}
    \item The optimal interest rate should be determined when marginal cost rate equals to marginal revenue rate
    \item Margin cost rate decrease at first and increases as the rises of interest rate
    \item Higher interest rate should always be taken therefore to appeal more deposits
    \item Amount of profit earned always decreases when interest rate increases
\end{choices}

\begin{correctanswer}
    A
\end{correctanswer}

\begin{explanation}
    \textbf{A选项正确。}
    这个说法正确：
    \begin{itemize}
        \item 最优利率在边际成本率等于边际收益率时确定
        \item 此时利润达到最大化
        \item 这是存款定价的基本原理
    \end{itemize}

    \textbf{B选项错误。}
    这个说法不准确：
    \begin{itemize}
        \item 边际成本率随利率上升而上升
        \item 不会先下降后上升
    \end{itemize}

    \textbf{C选项错误。}
    这个说法不准确：
    \begin{itemize}
        \item 高利率意味着高成本
        \item 不是利率越高越好
    \end{itemize}

    \textbf{D选项错误。}
    这个说法不准确：
    \begin{itemize}
        \item 利润先增加后减少
        \item 不是一直减少
    \end{itemize}
\end{explanation}

\checklist{存款利率定价原则}



\begin{question}
    Darcy is estimating the total financing cost included deposit and nondeposit funds. By reviewing the bank's past expense, he concluded that the following table:

    \begin{table}[htbp]
        \centering
        \begin{tabular}{|l|l|}
            \hline
            Total interest expense & \$50 million \\
            \hline
            Total operating expense & \$10 million \\
            \hline
            Shareholder's investment & \$150 million \\
            \hline
            Total funds raised & \$800 million \\
            \hline
            All earning assets & \$500 million \\
            \hline
            Required return of rate of shareholder (before tax) & 10\% \\
            \hline
            Tax rate & 35\% \\
            \hline
        \end{tabular}
        \caption{主要财务数据表}
    \end{table}

    Which of the following statement is NOT correct?
\end{question}

\begin{choices}
    \item The weighted average of interest expense is 6.25\%
    \item The break-even cost rate is 12\%
    \item The weighted average overall cost of capital is 15\%
    \item The least return of rate the bank should earn is 16.615\%
\end{choices}

\begin{correctanswer}
    D
\end{correctanswer}

\begin{explanation}
    \textbf{A选项错误。}
    这个说法不准确：
    \begin{itemize}
        \item 该说法正确
        \item 加权平均利息成本 = 50/800 = 6.25\%
    \end{itemize}

    \textbf{B选项错误。}
    这个说法不准确：
    \begin{itemize}
        \item 该说法正确
        \item 盈亏平衡成本率 = (50 + 10)/500 = 12\%
    \end{itemize}

    \textbf{C选项错误。}
    这个说法不准确：
    \begin{itemize}
        \item 该说法正确
        \item 加权平均总资本成本 = 12\% + 10\% × (150/500) = 15\%
    \end{itemize}

    \textbf{D选项正确。}
    这个说法正确：
    \begin{itemize}
        \item 该说法错误
        \item 最低回报率应为15\%
        \item 16.615\%的计算结果错误
    \end{itemize}
\end{explanation}

\checklist{融资成本计算}



\begin{question}
    A manager at a financial institution is assessing the cash flow requirements of the institution's largest customers using the available funds gap (AFG) model. The manager obtains the following current and projected inflows and outflows of funds:

    \begin{table}[htbp]
        \centering
        \begin{tabular}{|l|l|}
        \hline
        Deposits received today & USD 60 million \\
        \hline
        Approved new loan requests & USD 57 million \\
        \hline
        Expected drawdown on corporate client's credit lines & USD 16 million \\
        \hline
        Deposit inflows expected within the next week & USD 12 million \\
        \hline
        10-year Treasury securities expected to be purchased this week & USD 35 million \\
        \hline
        Other customer funds available today & USD 45 million \\
        \hline
        \end{tabular}
    \end{table}

    Assuming all else is held constant, if the institution's objective is to maintain an AFG of zero, what is the maximum amount of unexpected customer withdrawals it can meet within the next week?
\end{question}

\begin{choices}
    \item USD 0 million
    \item USD 3 million
    \item USD 9 million
    \item USD 15 million
\end{choices}

\begin{correctanswer}
    C
\end{correctanswer}

\begin{explanation}
    \textbf{A选项错误。}
    这个说法不准确：
    \begin{itemize}
        \item 计算结果错误
        \item 低估了可应对的提款金额
    \end{itemize}

    \textbf{B选项错误。}
    这个说法不准确：
    \begin{itemize}
        \item 计算结果错误
        \item 低估了可应对的提款金额
    \end{itemize}

    \textbf{C选项正确。}
    这个说法正确：
    \begin{itemize}
        \item 可用资金缺口(AFG)计算：
        \item 资金流入 = 60 + 12 + 45 = 117百万
        \item 资金流出 = 57 + 16 + 35 = 108百万
        \item 剩余资金 = 117 - 108 = 9百万
        \item 可应对的意外提款 = 9百万
    \end{itemize}

    \textbf{D选项错误。}
    这个说法不准确：
    \begin{itemize}
        \item 计算结果错误
        \item 高估了可应对的提款金额
    \end{itemize}
\end{explanation}

\checklist{可用资金缺口计算}




\newpage




\section{考点2:回购协议和融资}
\setcounter{question}{0}


\begin{question}
    In a presentation to management, a bond trader makes the following statements about repo collateral:
    \begin{enumerate}
        \item Statement 1: The difference between the federal funds rate and the general collateral rate is the special spread.
        \item Statement 2: During times of financial crisis, the spread between the federal funds rate and the general collateral rate widens.
    \end{enumerate}
    Which of the trader's statements are accurate?
\end{question}

\begin{choices}
    \item Both statements are incorrect
    \item Only statement 1 is correct
    \item Only statement 2 is correct
    \item Both statement are correct
\end{choices}

\begin{correctanswer}
    C
\end{correctanswer}

\begin{explanation}
    \textbf{A选项错误。}
    这个说法不准确：
    \begin{itemize}
        \item 第二个陈述是正确的
        \item 金融危机期间利差确实扩大
    \end{itemize}

    \textbf{B选项错误。}
    这个说法不准确：
    \begin{itemize}
        \item 第一个陈述是错误的
        \item 联邦基金利率和GC利率的差额是fed funds-GC利差
        \item 特殊利差是GC利率和特定证券特殊利率的差额
    \end{itemize}

    \textbf{C选项正确。}
    这个说法正确：
    \begin{itemize}
        \item 第一个陈述错误
        \item 第二个陈述正确
        \item 金融危机期间，借出国债意愿下降
        \item 导致GC利率下降，利差扩大
    \end{itemize}

    \textbf{D选项错误。}
    这个说法不准确：
    \begin{itemize}
        \item 第一个陈述是错误的
        \item 混淆了特殊利差和fed funds-GC利差
    \end{itemize}
\end{explanation}

\checklist{回购担保品利差分析}


\begin{question}
    In contrast to general collateral (GC) repo rates, which of the following is TRUE about special repo rates?
\end{question}

\begin{choices}
    \item Special rates are typically less than general collateral rates
    \item If the counterparty's primary motivation is to lend cash rather than borrow a security, the special rate applies
    \item Special rates are well-suited to repo investors who are looking to obtain the highest rate for the collateral they are willing to accept
    \item The most commonly cited special rates are for overnight repos where any U.S. Treasury collateral is acceptable
\end{choices}

\begin{correctanswer}
    A
\end{correctanswer}

\begin{explanation}
    \textbf{A选项正确。}
    这个说法正确：
    \begin{itemize}
        \item 特殊利率通常低于一般抵押利率
        \item 这是因为特殊证券的需求更高
        \item 借款人愿意接受较低利率以获得特定证券
    \end{itemize}

    \textbf{B选项错误。}
    这个说法不准确：
    \begin{itemize}
        \item 特殊利率适用于需要特定证券的情况
        \item 与主要动机是借出资金还是借入证券无关
    \end{itemize}

    \textbf{C选项错误。}
    这个说法不准确：
    \begin{itemize}
        \item 特殊利率通常较低
        \item 不适合追求最高利率的投资者
    \end{itemize}

    \textbf{D选项错误。}
    这个说法不准确：
    \begin{itemize}
        \item 特殊利率适用于特定证券
        \item 不是任何国债都适用
    \end{itemize}
\end{explanation}

\checklist{特殊回购利率特征}


\begin{question}
    How many of the following statements is most likely correct?
    \begin{enumerate}
        \item I. The special spreads usually derives from the difference between the general collateral rate and the repo rate if the collateral securities are on-the-run (OTR) T-bonds.
        \item II. An important reason that special trade usually uses on-the-run (OTR) bonds is the brilliant liquidity involved in new-issued T-bonds, and the liquidity has been cherished by both investors who hold long positions and short sellers.
        \item III. Special spread would usually become largest immediately after auctions and smallest before auctions.
    \end{enumerate}
\end{question}

\begin{choices}
    \item Zero
    \item One
    \item Two
    \item Three
\end{choices}

\begin{correctanswer}
    C
\end{correctanswer}

\begin{explanation}
    \textbf{A选项错误。}
    这个说法不准确：
    \begin{itemize}
        \item 有两个陈述是正确的
        \item 低估了正确陈述的数量
    \end{itemize}

    \textbf{B选项错误。}
    这个说法不准确：
    \begin{itemize}
        \item 有两个陈述是正确的
        \item 低估了正确陈述的数量
    \end{itemize}

    \textbf{C选项正确。}
    这个说法正确：
    \begin{itemize}
        \item I正确：特殊利差确实来自GC利率和回购利率的差额
        \item II正确：OTR债券的流动性确实受到多空双方的重视
        \item III错误：利差在拍卖后最小，拍卖前最大
        \item 因为拍卖后新OTR证券的供应增加，抑制了特殊利差
    \end{itemize}

    \textbf{D选项错误。}
    这个说法不准确：
    \begin{itemize}
        \item 只有两个陈述是正确的
        \item 高估了正确陈述的数量
    \end{itemize}
\end{explanation}

\checklist{特殊利差特征分析}



\begin{question}
    In regard to special spreads, each of the following is true EXCEPT which is false?
\end{question}

\begin{choices}
    \item On-the-run (OTR) issues tend to trade "more special" than off-the-run (OFR) issues, where "more special" refers to special spreads that are larger
    \item The special spread equals the general collateral (GC) repo rate minus the special collateral repo rate
    \item On-the-run special spreads peak immediately after an auction, and tend to decrease over the cycle, reaching their lowest level immediately before the next auction
    \item Special spreads tend to be volatile on a daily basis and special spreads can be quite large (e.g., hundreds of basis points)
\end{choices}

\begin{correctanswer}
    C
\end{correctanswer}

\begin{explanation}
    \textbf{A选项错误。}
    这个说法不准确：
    \begin{itemize}
        \item 该说法正确
        \item OTR债券确实比OFR债券有更大的特殊利差
    \end{itemize}

    \textbf{B选项错误。}
    这个说法不准确：
    \begin{itemize}
        \item 该说法正确
        \item 特殊利差确实等于GC回购利率减去特殊抵押品回购利率
    \end{itemize}

    \textbf{C选项正确。}
    这个说法正确：
    \begin{itemize}
        \item 该说法错误
        \item 实际情况相反：
        \item 拍卖后利差最小
        \item 周期内逐渐增加
        \item 下次拍卖前达到峰值
    \end{itemize}

    \textbf{D选项错误。}
    这个说法不准确：
    \begin{itemize}
        \item 该说法正确
        \item 特殊利差确实波动较大
        \item 可能达到数百个基点
    \end{itemize}
\end{explanation}

\checklist{特殊利差特征分析}


\begin{question}
    Rankcon is a repo investor in the money market mutual fund industry with the exclusive motive of cash management. Given this motivation, Rankcon understandably employs each of the following criteria (or preference) with respect to its repo investments EXCEPT which is the LEAST likely preference?
\end{question}

\begin{choices}
    \item Rankcon insists on a sufficient collateral haircut
    \item Rankcon only accept securities with the highest credit quality; e.g., debt of government-sponsored entities (GSEs)
    \item Rankcon refuses to accept general collateral and instead insists on particular securities with delineated asset classes
    \item Rankcon places a premium on liquidity so tends to lend overnight rather than for term; or, if it wants to lend cash for an extended period, engages in an open repo
\end{choices}

\begin{correctanswer}
    C
\end{correctanswer}

\begin{explanation}
    \textbf{A选项错误。}
    这个说法不准确：
    \begin{itemize}
        \item 该说法正确
        \item 要求足够的抵押品折扣是合理的
        \item 可以降低对手方违约风险
    \end{itemize}

    \textbf{B选项错误。}
    这个说法不准确：
    \begin{itemize}
        \item 该说法正确
        \item 只接受最高信用质量的证券是合理的
        \item 符合现金管理的安全需求
    \end{itemize}

    \textbf{C选项正确。}
    这个说法正确：
    \begin{itemize}
        \item 该说法错误
        \item 回购投资者通常不要求特定债券
        \item 只关注抵押品类别而非具体证券
        \item 一般接受普通抵押品
    \end{itemize}

    \textbf{D选项错误。}
    这个说法不准确：
    \begin{itemize}
        \item 该说法正确
        \item 重视流动性是合理的
        \item 倾向于隔夜或开放式回购
    \end{itemize}
\end{explanation}

\checklist{回购投资偏好分析}



\newpage




\section{考点3:金融服务管理中的投资职能}
\setcounter{question}{0}

\begin{question}
    The main business lines of large dealer banks are: (1) securities dealing, underwriting, and trading; (2) over-the-counter derivatives; (3) prime brokerage and asset management. About these main business lines, each of the following statements is true except which is false?
\end{question}

\begin{choices}
    \item A tri-party repo agreement is used to mitigate counterparty risk and/or reduce burden of collateral management
    \item For most OTC derivatives, one of the counterparties is a dealer who typically lays off much of the risk of its client-initiated trades by running a matched book
    \item With respect to OTC derivatives, the best measure of their systemic risk is the total market value of all outstanding contracts; further, master swap agreements do not reduce their actual risks
    \item Some large dealers are Prime Brokers who provide clients a range of services including management of securities holdings, clearing, cash-management, securities lending, financing, reporting, and tax accounting
\end{choices}

\begin{correctanswer}
    C
\end{correctanswer}

\begin{explanation}
    \textbf{A选项错误。}
    这个说法不准确：
    \begin{itemize}
        \item 该说法正确
        \item 三方回购确实用于降低对手方风险
        \item 减轻抵押品管理负担
    \end{itemize}

    \textbf{B选项错误。}
    这个说法不准确：
    \begin{itemize}
        \item 该说法正确
        \item 交易商确实通过匹配账簿对冲风险
        \item 这是OTC衍生品交易的标准做法
    \end{itemize}

    \textbf{C选项正确。}
    这个说法正确：
    \begin{itemize}
        \item 该说法错误
        \item 主协议通过净额结算降低风险
        \item 标准化抵押品安排也降低风险
        \item 所有衍生品合约的总市值为零
    \end{itemize}

    \textbf{D选项错误。}
    这个说法不准确：
    \begin{itemize}
        \item 该说法正确
        \item 主要经纪商确实提供全面的服务
        \item 包括证券管理、清算、现金管理等
    \end{itemize}
\end{explanation}

\checklist{交易商银行业务分析}

\begin{question}
    In recent year, large dealer banks financed significant fractions of their assets using short-term, often overnight, repurchase (repo) agreements in which creditors held bank securities as collateral against default losses. The table below shows the quarter-end financing of four broker-dealer banks. All values are in USD billions.

    \begin{table}[htbp]
        \centering
        \begin{tabular}{|l|c|c|c|c|}
            \hline
            & \textbf{Bank A} & \textbf{Bank B} & \textbf{Bank C} & \textbf{Bank D} \\
            \hline
            Financial Instruments Owned & 823 & 629 & 723 & 382 \\
            \hline
            Pledged as collateral      & 272 & 289 & 380 & 155 \\
            \hline
        \end{tabular}
    \end{table}
    If repo creditors become nervous about a bank's solvency, which bank is least vulnerable to a liquidity crisis?
\end{question}

\begin{choices}
    \item Bank A
    \item Bank B
    \item Bank C
    \item Bank D
\end{choices}

\begin{correctanswer}
    A
\end{correctanswer}

\begin{explanation}
    \textbf{A选项正确。}
    这个说法正确：
    \begin{itemize}
        \item 银行A的抵押品比例最低
        \item 272/823 = 33\%
        \item 对短期回购融资依赖最小
        \item 流动性风险最低
    \end{itemize}

    \textbf{B选项错误。}
    这个说法不准确：
    \begin{itemize}
        \item 银行B的抵押品比例为46\%
        \item 289/629 = 46\%
        \item 对短期回购融资依赖较大
    \end{itemize}

    \textbf{C选项错误。}
    这个说法不准确：
    \begin{itemize}
        \item 银行C的抵押品比例为53\%
        \item 380/723 = 53\%
        \item 对短期回购融资依赖最大
    \end{itemize}

    \textbf{D选项错误。}
    这个说法不准确：
    \begin{itemize}
        \item 银行D的抵押品比例为41\%
        \item 155/382 = 41\%
        \item 对短期回购融资依赖较大
    \end{itemize}
\end{explanation}

\checklist{回购融资风险分析}


\begin{question}
    The Rosenfeld Investment Firm wants to invest in bond portfolios. Based on its multi-year fabulous expertise in bond investment, it can choose from among the following maturity strategies: Ladder Policy, Front-end Load Maturity Policy, Back-end Load Maturity Policy, Barbell Strategy, or Rate Expectations Approach. The firm's goal for the portfolio is neither to maximize income nor to seek to maximize the upside potential for earnings. Instead, the primary goal is to use the portfolio as a source of liquidity. For that purpose, which of the following strategy is the most suitable?
\end{question}

\begin{choices}
    \item Barbell strategy
    \item Front-end Load Maturity Policy
    \item Back-end Load Maturity Policy
    \item Ladder policy
\end{choices}

\begin{correctanswer}
    B
\end{correctanswer}

\begin{explanation}
    \textbf{A选项错误。}
    这个说法不准确：
    \begin{itemize}
        \item 杠铃策略不适合流动性需求
        \item 主要投资于短期和长期债券
        \item 中间期限债券较少
    \end{itemize}

    \textbf{B选项正确。}
    这个说法正确：
    \begin{itemize}
        \item 前端加载期限政策最适合
        \item 主要购买短期证券
        \item 短期到期提供流动性
        \item 符合公司的流动性目标
    \end{itemize}

    \textbf{C选项错误。}
    这个说法不准确：
    \begin{itemize}
        \item 后端加载期限政策不适合
        \item 主要投资于长期债券
        \item 流动性较差
    \end{itemize}

    \textbf{D选项错误。}
    这个说法不准确：
    \begin{itemize}
        \item 阶梯策略不适合
        \item 投资于多个期限
        \item 流动性不如前端加载策略
    \end{itemize}
\end{explanation}

\checklist{债券投资策略分析}




\newpage




\section{考点4:非流动资产}
\setcounter{question}{0}


\begin{question}
    One of the board members has suggested that the CIO look at several hedge funds that have reported strong performance in recent years. The CRO is familiar with many of these funds and is aware that several invest heavily in illiquid assets and that this may cause standard risk measures based on daily returns to give a misleading picture of their risk. Which of following statements about daily risk measures applied to hedge funds that invest in illiquid assets is correct?
\end{question}

\begin{choices}
    \item Correlation with other investments will be artificially lowered, giving the appearance of low systematic risk
    \item Infrequent trading reduces the smoothing effects from mark-to-market valuation, giving the appearance of high volatility
    \item Returns of illiquid assets tend to exhibit negative serial correlation, leading to higher long-term volatility
    \item All else being equal, illiquid assets tend to have lower Sharpe ratios, causing the incorrect appearance of an illiquid premium
\end{choices}

\begin{correctanswer}
    A
\end{correctanswer}

\begin{explanation}
    \textbf{A选项正确。}
    这个说法正确：
    \begin{itemize}
        \item 非流动性资产交易不频繁
        \item 导致与其他投资的相关性被低估
        \item 系统性风险被低估
    \end{itemize}

    \textbf{B选项错误。}
    这个说法不准确：
    \begin{itemize}
        \item 不频繁交易使数据更平滑
        \item 导致波动率被低估
        \item 不是高估波动率
    \end{itemize}

    \textbf{C选项错误。}
    这个说法不准确：
    \begin{itemize}
        \item 非流动性资产收益通常呈现正序列相关
        \item 不是负序列相关
        \item 长期波动率被低估
    \end{itemize}

    \textbf{D选项错误。}
    这个说法不准确：
    \begin{itemize}
        \item 非流动性资产通常有更高的夏普比率
        \item 不是更低的夏普比率
        \item 流动性溢价被高估
    \end{itemize}
\end{explanation}

\checklist{非流动性资产风险度量}


\begin{question}
    The treasurer of ABC, a large industrial firm, is considering issuing a 7-year bond to fund a planned expansion. In order for the debt issuance to be successful, the treasurer feels that the bond must be attractive to investors in terms of potential return and liquidity. To get a sense of the current market and how investors might receive the new debt issuance, the treasurer asks the risk department to conduct an analysis of the historical trading activity and performance of bonds issued by firms similar to ABC. The bonds range in their initial maturity from 6 to 8 years but vary significantly in liquidity with some bonds traded several times a week and others that may only trade a few times a year. When analyzing the information, the risk department should be concerned that an illiquid bond, when compared to a liquid bond, could exhibit:
\end{question}

\begin{choices}
    \item Lower returns and bid-ask spreads
    \item Higher autocorrelation of returns
    \item A lower calculated Sharpe ratio
    \item A higher beta to the market
\end{choices}

\begin{correctanswer}
    B
\end{correctanswer}

\begin{explanation}
    \textbf{A选项错误。}
    这个说法不准确：
    \begin{itemize}
        \item 非流动性债券通常有更高的买卖价差
        \item 不是更低的买卖价差
        \item 回报率可能更高
    \end{itemize}

    \textbf{B选项正确。}
    这个说法正确：
    \begin{itemize}
        \item 非流动性债券收益呈现更高的自相关性
        \item 市场采用收入平滑方式报告收益
        \item 当前数据与历史数据高度相关
    \end{itemize}

    \textbf{C选项错误。}
    这个说法不准确：
    \begin{itemize}
        \item 非流动性资产通常有更高的夏普比率
        \item 不是更低的夏普比率
        \item 波动率被低估
    \end{itemize}

    \textbf{D选项错误。}
    这个说法不准确：
    \begin{itemize}
        \item 非流动性资产通常有更低的市场贝塔
        \item 不是更高的市场贝塔
        \item 与市场相关性较低
    \end{itemize}
\end{explanation}

\checklist{非流动性债券特征分析}


\begin{question}
    A risk analyst in a hedge fund is gauging the liquidity risk exposure of a portfolio which allocates on illiquid assets such as private equity, convertible bond, treasure assets, etc. Which one of following statements about risk and return biases of illiquid assets is incorrect?
\end{question}

\begin{choices}
    \item The risk analyst found a significant first-order autocorrelation coefficient of 0.5 for the quarterly historical returns. This can be seen as an indicator of historical return smoothing
    \item For an illiquid asset, the correlation between asset return and market return will be artificially biased upward, giving the appearance of high systematic risk
    \item If there is survivorship bias in this portfolio, the returns tend to be biased upward
    \item The price impact of large trades can be seen as an illiquidity variable to measure the illiquidity in equity market
\end{choices}

\begin{correctanswer}
    B
\end{correctanswer}

\begin{explanation}
    \textbf{A选项错误。}
    这个说法不准确：
    \begin{itemize}
        \item 该说法正确
        \item 显著的自相关系数表明收益平滑
        \item 这是非流动性资产的典型特征
    \end{itemize}

    \textbf{B选项正确。}
    这个说法正确：
    \begin{itemize}
        \item 该说法错误
        \item 非流动性资产与市场的相关性被低估
        \item 不是被高估
        \item 系统性风险被低估
    \end{itemize}

    \textbf{C选项错误。}
    这个说法不准确：
    \begin{itemize}
        \item 该说法正确
        \item 幸存者偏差确实导致收益高估
        \item 只保留表现好的资产数据
    \end{itemize}

    \textbf{D选项错误。}
    这个说法不准确：
    \begin{itemize}
        \item 该说法正确
        \item 大额交易的市场影响是流动性指标
        \item 影响越小流动性越好
    \end{itemize}
\end{explanation}

\checklist{非流动性资产偏差分析}


\begin{question}
    Consider investors in the following five asset classes:
    \begin{enumerate}
        \item Leveraged buyout (LBO) investors
        \item Merger arbitrage hedge funds
        \item Mortgage loan investors
        \item Convertible bond investors
        \item Statistical arbitrage investors
    \end{enumerate}
    Which of the above is (are) materially or meaningfully exposed to systematic funding liquidity risk?
\end{question}

\begin{choices}
    \item None of the above
    \item I. and II. Only
    \item III. and IV. only
    \item All of the above
\end{choices}

\begin{correctanswer}
    D
\end{correctanswer}

\begin{explanation}
    \textbf{A选项错误。}
    这个说法不准确：
    \begin{itemize}
        \item 所有资产类别都面临系统性融资流动性风险
        \item 低估了风险暴露
    \end{itemize}

    \textbf{B选项错误。}
    这个说法不准确：
    \begin{itemize}
        \item 不仅LBO和合并套利面临风险
        \item 其他资产类别也面临风险
        \item 低估了风险暴露
    \end{itemize}

    \textbf{C选项错误。}
    这个说法不准确：
    \begin{itemize}
        \item 不仅抵押贷款和可转债面临风险
        \item 其他资产类别也面临风险
        \item 低估了风险暴露
    \end{itemize}

    \textbf{D选项正确。}
    这个说法正确：
    \begin{itemize}
        \item LBO投资者面临融资失败风险
        \item 合并套利基金受融资环境影响
        \item 抵押贷款投资者面临再融资风险
        \item 可转债投资者受市场流动性影响
        \item 统计套利策略依赖融资环境
    \end{itemize}
\end{explanation}

\checklist{系统性融资流动性风险分析}


\begin{question}
    A fund manager is presenting to a group of analysts on the topic of asset illiquidity. The manager highlights the sources of illiquidity in the presentation and describes the relationship between market imperfections and illiquidity. Which of the following is a correct statement to include in the presentation?
\end{question}

\begin{choices}
    \item Insufficient capital and differences in investor expertise can result in illiquid asset markets
    \item Asset illiquidity is independent of investors' access to credit
    \item Transactions, whether large or small in size, will have an equal impact on asset prices
    \item Symmetric information among a large group of investors can cause markets to freeze and break down
\end{choices}

\begin{correctanswer}
    A
\end{correctanswer}

\begin{explanation}
    \textbf{A选项正确。}
    这个说法正确：
    \begin{itemize}
        \item 资本不足是市场非流动性的重要来源
        \item 投资者专业知识和经验差异影响市场流动性
        \item 只有具备足够资本和专业的投资者才能参与
    \end{itemize}

    \textbf{B选项错误。}
    这个说法不准确：
    \begin{itemize}
        \item 融资限制是非流动性的重要来源
        \item 信贷可得性直接影响市场流动性
        \item 无法获得信贷的投资者难以交易
    \end{itemize}

    \textbf{C选项错误。}
    这个说法不准确：
    \begin{itemize}
        \item 大额交易对价格影响显著
        \item 小额交易影响较小
        \item 交易规模影响不等
    \end{itemize}

    \textbf{D选项错误。}
    这个说法不准确：
    \begin{itemize}
        \item 信息不对称导致市场冻结
        \item 不是信息对称
        \item 投资者担心被剥削而不愿交易
    \end{itemize}
\end{explanation}

\checklist{资产非流动性来源分析}


\newpage


\section{考点5:针对利率变化的风险管理：资产负债管理和期限技术}
\setcounter{question}{0}
\begin{question}
    Changing interest rates also changes the market value of assets and liabilities, thereby changing each financial institution's net worth. Shareholders may pay more attention to the management of net worth, which of the following statements about duration gap management is least correct?
\end{question}

\begin{choices}
    \item The larger the leverage-adjusted duration gap, the more sensitive will be the net worth of a financial institution to a change in interest rates
    \item A financial firm with longer-duration assets than liabilities will suffer a greater decline in net worth when market interest rates rise
    \item Net worth is more sensitive to interest rate when rate is at high level than at low due to the existence of convexity
    \item In practice, duration gap management considers only the liner relationship between the change of interest rate and the asset/liability price, which is a limitation
\end{choices}

\begin{correctanswer}
    C
\end{correctanswer}

\begin{explanation}
    \textbf{A选项错误。}
    这个说法不准确：
    \begin{itemize}
        \item 该说法正确
        \item 杠杆调整后的久期缺口越大
        \item 净资产对利率变化越敏感
    \end{itemize}

    \textbf{B选项错误。}
    这个说法不准确：
    \begin{itemize}
        \item 该说法正确
        \item 资产久期大于负债久期时
        \item 利率上升会导致净资产下降
    \end{itemize}

    \textbf{C选项正确。}
    这个说法正确：
    \begin{itemize}
        \item 该说法错误
        \item 由于凸性存在
        \item 低利率时净资产对利率更敏感
        \item 不是高利率时更敏感
    \end{itemize}

    \textbf{D选项错误。}
    这个说法不准确：
    \begin{itemize}
        \item 该说法正确
        \item 久期缺口管理只考虑线性关系
        \item 这是一个局限性
    \end{itemize}
\end{explanation}

\checklist{久期缺口管理分析}


\begin{question}
    Asset-liability committee (ALCO), meets regularly to manage the financial firm's interest rate risk and other risk exposure, and estimates the firm's risk exposure. As a member of ALCO, Jason finds that the bank has negative interest-sensitive gap, which of the following statements is correct?
\end{question}

\begin{choices}
    \item The bank is asset sensitive
    \item If interest rates rise, bank's net interest margin will decrease
    \item If interest rates rise, bank's net interest margin will increase
    \item The Interest Sensitivity Ratio (ISR) is more than 1
\end{choices}

\begin{correctanswer}
    B
\end{correctanswer}

\begin{explanation}
    \textbf{A选项错误。}
    这个说法不准确：
    \begin{itemize}
        \item 负利率敏感缺口表示负债敏感
        \item 不是资产敏感
        \item 利率敏感负债多于利率敏感资产
    \end{itemize}

    \textbf{B选项正确。}
    这个说法正确：
    \begin{itemize}
        \item 负利率敏感缺口表示负债敏感
        \item 利率上升时负债利息支出增加更多
        \item 资产利息收入增加较少
        \item 导致净息差下降
    \end{itemize}

    \textbf{C选项错误。}
    这个说法不准确：
    \begin{itemize}
        \item 利率上升时净息差会下降
        \item 不是增加
        \item 负债敏感型银行的特点
    \end{itemize}

    \textbf{D选项错误。}
    这个说法不准确：
    \begin{itemize}
        \item 负利率敏感缺口时
        \item 利率敏感比率小于1
        \item 不是大于1
    \end{itemize}
\end{explanation}

\checklist{利率敏感缺口分析}


\begin{question}
    No financial manager can completely avoid one of the toughest and potentially most damaging forms of risk that all financial institutions must face—interest rate risk. An analyst is evaluating four banks interest risk management, the information is as follows:

    \begin{table}[htbp]
        \centering
        \begin{tabular}{|l|c|c|}
            \hline
            & \textbf{Bank A} & \textbf{Bank B} \\
            \hline
            Interest-sensitive gap & \$50 million & \$-30 million \\
            \hline
        \end{tabular}
    \end{table}

    Which of the following statements are appropriate strategies these banks can adopt?
    \begin{enumerate}
        \item I. Bank A should increase investing in interest-sensitive assets if rate increases 1\%
        \item II. Bank B should decrease its interest-sensitive liabilities if rate decreases 2\%
        \item III. Bank C should purchase more assets with longer maturity and decreasing liabilities with longer maturity if rate falls from 4\% to 2\%
        \item IV. Bank D should decrease duration of assets and increase duration of liabilities if rate raises from 2\% to 4\%
    \end{enumerate}
\end{question}

\begin{choices}
    \item I, III \& IV
    \item II \& IV
    \item I, II \& III
    \item I, II, III \& IV
\end{choices}

\begin{correctanswer}
    A
\end{correctanswer}

\begin{explanation}
    \textbf{A选项正确。}
    这个说法正确：
    \begin{itemize}
        \item I正确：正利率敏感缺口时，增加利率敏感资产可增加净利息收入
        \item III正确：利率下降时，增加资产久期和减少负债久期可最大化净资产
        \item IV正确：利率上升时，减少资产久期和增加负债久期可最大化净资产
    \end{itemize}

    \textbf{B选项错误。}
    这个说法不准确：
    \begin{itemize}
        \item II错误：负利率敏感缺口时，利率下降应增加利率敏感负债
        \item 遗漏了I和III的正确策略
    \end{itemize}

    \textbf{C选项错误。}
    这个说法不准确：
    \begin{itemize}
        \item II错误：负利率敏感缺口时，利率下降应增加利率敏感负债
        \item 遗漏了IV的正确策略
    \end{itemize}

    \textbf{D选项错误。}
    这个说法不准确：
    \begin{itemize}
        \item II错误：负利率敏感缺口时，利率下降应增加利率敏感负债
        \item 包含了错误的策略
    \end{itemize}
\end{explanation}

\checklist{利率风险管理策略分析}


\begin{question}
    A risk consultant is reviewing the credit function at Bank WDI. The bank recently shifted its lending business from using the originate-to-hold model to using the originate-to-distribute model. Under the originate-to-distribute model, Bank WDI divides its loans into core loans that it holds until maturity, and noncore loans that it hedges by means of credit derivatives or sells by means of loan securitization. Which of the following statements would the consultant be correct to make?
\end{question}

\begin{choices}
    \item The non-core loans are managed by the originating business units and the core loans are transfer priced to the credit portfolio management group
    \item When economic capital is allocated to the core loans, the allocation is based on the weighted average of the loan exposures, with the weights determined by the size of the core loans
    \item When any core loan is originated, a primary consideration is that its spread should produce a risk-adjusted return on capital that is greater than the hurdle rate of the bank
    \item The spread charged to a noncore loan at origination is directly related to the loan's credit rating and maturity
\end{choices}

\begin{correctanswer}
    C
\end{correctanswer}

\begin{explanation}
    \textbf{A选项错误。}
    这个说法不准确：
    \begin{itemize}
        \item 核心贷款由发起业务部门管理
        \item 非核心贷款转移定价给信贷组合管理组
        \item 说法与实际情况相反
    \end{itemize}

    \textbf{B选项错误。}
    这个说法不准确：
    \begin{itemize}
        \item 经济资本分配基于风险加权
        \item 不是基于贷款规模加权
        \item 需要考虑信用风险
    \end{itemize}

    \textbf{C选项正确。}
    这个说法正确：
    \begin{itemize}
        \item 核心贷款发起时需考虑风险调整后收益
        \item 收益需超过银行的最低要求回报率
        \item 这是合理的风险管理要求
    \end{itemize}

    \textbf{D选项错误。}
    这个说法不准确：
    \begin{itemize}
        \item 非核心贷款利差与市场条件相关
        \item 不是直接与信用评级和期限相关
        \item 需要考虑市场流动性
    \end{itemize}
\end{explanation}

\checklist{贷款管理模式分析}



\begin{question}

\end{question}

\begin{choices}
    \item 
\end{choices}

\begin{correctanswer}
    B
\end{correctanswer}

\begin{explanation}
    \textbf{A选项错误。}
    \begin{itemize}
        \item 
    \end{itemize}
    \textbf{B选项错误。}
    \begin{itemize}
        \item 
    \end{itemize}
    \textbf{C选项错误。}
    \begin{itemize}
        \item 
    \end{itemize}
    \textbf{D选项错误。}
    \begin{itemize}
        \item 
    \end{itemize}
\end{explanation}
\checklist{}



\begin{question}

\end{question}

\begin{choices}
    \item 
\end{choices}

\begin{correctanswer}
    B
\end{correctanswer}

\begin{explanation}
    \textbf{A选项错误。}
    \begin{itemize}
        \item 
    \end{itemize}
    \textbf{B选项错误。}
    \begin{itemize}
        \item 
    \end{itemize}
    \textbf{C选项错误。}
    \begin{itemize}
        \item 
    \end{itemize}
    \textbf{D选项错误。}
    \begin{itemize}
        \item 
    \end{itemize}
\end{explanation}
\checklist{}



\begin{question}

\end{question}

\begin{choices}
    \item 
\end{choices}

\begin{correctanswer}
    B
\end{correctanswer}

\begin{explanation}
    \textbf{A选项错误。}
    \begin{itemize}
        \item 
    \end{itemize}
    \textbf{B选项错误。}
    \begin{itemize}
        \item 
    \end{itemize}
    \textbf{C选项错误。}
    \begin{itemize}
        \item 
    \end{itemize}
    \textbf{D选项错误。}
    \begin{itemize}
        \item 
    \end{itemize}
\end{explanation}
\checklist{}



\begin{question}

\end{question}

\begin{choices}
    \item 
\end{choices}

\begin{correctanswer}
    B
\end{correctanswer}

\begin{explanation}
    \textbf{A选项错误。}
    \begin{itemize}
        \item 
    \end{itemize}
    \textbf{B选项错误。}
    \begin{itemize}
        \item 
    \end{itemize}
    \textbf{C选项错误。}
    \begin{itemize}
        \item 
    \end{itemize}
    \textbf{D选项错误。}
    \begin{itemize}
        \item 
    \end{itemize}
\end{explanation}
\checklist{}



\begin{question}

\end{question}

\begin{choices}
    \item 
\end{choices}

\begin{correctanswer}
    B
\end{correctanswer}

\begin{explanation}
    \textbf{A选项错误。}
    \begin{itemize}
        \item 
    \end{itemize}
    \textbf{B选项错误。}
    \begin{itemize}
        \item 
    \end{itemize}
    \textbf{C选项错误。}
    \begin{itemize}
        \item 
    \end{itemize}
    \textbf{D选项错误。}
    \begin{itemize}
        \item 
    \end{itemize}
\end{explanation}
\checklist{}


\newpage


\end{document}